%% chapters/01-introduction.tex

\section{Introduction}
\label{sec:intro}

Large language models (LLMs)—including the GPT family, Llama, Gemini, and Claude—have rapidly migrated from research prototypes to broadly deployed decision-support tools across high-stakes domains such as healthcare, finance, and law. Built on massive neural architectures and web-scale corpora, these systems exhibit strong capabilities in language understanding, generation, and multi-step reasoning. Their expanding societal integration, however, raises equally expanding risks: models can produce harmful or biased content, leak sensitive information, or fabricate plausible but false claims, with real-world consequences when used in critical workflows. We define trustworthiness as the reliability of an LLM in deployment—safe behavior, privacy protection, fairness, and factual accuracy. Ensuring these properties demands behavioral control at inference time, not just greater model capability. Developing control mechanisms that are both practical and scalable is therefore a primary scientific and engineering challenge for LLM deployment.

% The rapid proliferation and societal integration of Large Language Models (LLMs) such as the GPT-series, Llama, Gemini, and Claude represent a paradigm shift in artificial intelligence. These models, built upon vast neural architectures and trained on petabytes of data, have demonstrated remarkable capabilities in natural language understanding, generation, and complex reasoning. Their deployment is no longer confined to academic research or niche applications; they are starting to be utilized as decision-support tools in critical societal systems, including healthcare, finance, and law. This deep integration into high-stakes environments means that model failures can have severe consequences. Issues such as generating harmful content, leaking private data, amplifying societal biases, or fabricating misinformation pose significant risks, collectively challenging the models' trustworthiness. As such, the ability to effectively guide and control model behavior to ensure its reliability and alignment with human values—in essence, to guarantee its trustworthiness—has emerged as a central scientific and engineering challenge.
% As this integration deepens, the ability to effectively guide and control model behavior to ensure its reliability and alignment with human values has emerged as a central scientific and engineering challenge.

The primary approach to this challenge has been training-time alignment: methods such as instruction fine-tuning and Reinforcement Learning from Human Feedback (RLHF) encode desired behaviors by updating model parameters. These techniques have been pivotal in turning raw pretrained models into helpful and safe assistants. However, they face inherent limits. Once trained, behavior is largely fixed and cannot adapt to new policies, users, or contexts without retraining; retraining itself is expensive; and for proprietary, API-only systems, it is often infeasible. These constraints motivate post-training control during deployment, shifting attention to methods that operate at or immediately around inference.

% The primary approach to this challenge has been training-time alignment, where techniques like instruction fine-tuning and Reinforcement Learning from Human Feedback (RLHF) embed desired behaviors by directly modifying the model's internal parameters through gradient-based optimization. These methods have been instrumental in transforming raw pre-trained models into helpful and safe assistants. However, they possess inherent limitations: they are static, as model behavior is largely fixed post-training; costly, requiring immense computational resources for retraining; and often inapplicable to proprietary, closed-source models accessible only via APIs. These constraints have consequently spurred research into alternative control strategies that can be applied dynamically after training, shifting the focus towards techniques that operate during inference.
% These constraints have catalyzed a growing shift in focus towards methods that operate during inference.

Historically, inference-time work has focused on computational efficiency—lower latency, smaller memory footprint, and reduced energy use. Building on these foundations, a complementary paradigm is emerging: repurposing inference-time techniques as a behavioral control layer that targets trustworthiness. Instead of optimizing throughput alone, these methods are explicitly crafted to safeguard safety, privacy, fairness, and factuality at deployment. We argue that such interventions augment, rather than replace, training-time alignment by enabling real-time, adaptive, and modular control of model behavior in practice.

% Historically, inference-time research primarily emphasized efficiency objectives such as reducing latency, memory footprint, or energy consumption. A new paradigm is now emerging: re-purposing inference-time methods as a behavioral control layer for trustworthiness. Rather than focusing solely on performance, these methods are increasingly developed to safeguard safety, privacy, fairness, factuality, and faithfulness during deployment. This survey explores this frontier, highlighting how inference-time interventions can complement training-time alignment by providing real-time, adaptive, and modular control over model behavior.

To structure our analysis, we focus on four core dimensions of trustworthiness, each representing a critical area of potential risk and societal impact:

\begin{itemize}
    \item \textbf{Safety:} Prevent the generation of harmful, toxic, biased, or malicious content; mitigate misuse for disinformation or dangerous activities; and defend against adversarial threats, including \emph{jailbreak attacks} that attempt to bypass safety policies and \emph{backdoor attacks} that embed hidden malicious behaviors.
    \item \textbf{Privacy:} Protect sensitive information by preventing memorization and leakage of personal data from pre-training corpora, and safeguard the confidentiality of user inputs during inference.
    \item \textbf{Fairness:} Avoid the perpetuation or amplification of societal biases, ensuring that outputs do not systematically disadvantage individuals or groups on the basis of protected attributes.
    \item \textbf{Factuality:} Produce outputs that are grounded in verifiable evidence and consistent with established world knowledge, thereby minimizing factually incorrect or hallucinated content.
    % \item \textbf{Faithfulness:} Ensuring outputs remain strictly grounded in and logically consistent with the provided context and user instructions, preventing hallucination or reasoning that deviates from the immediate task constraints. This is distinct from factuality, which measures alignment with external world knowledge.
\end{itemize}

Addressing these multidimensional risks requires techniques that are dynamic, lightweight, and compatible with frozen pretrained models. We define inference-time methods as post-training interventions that guide, constrain, or steer model behavior during or immediately before generation, without full retraining. These methods span a spectrum of intrusiveness. At one end, purely black-box approaches (e.g., prompting, guardrails, decoding) operate on inputs and outputs only. At the other end, lightweight but more intrusive interventions (e.g., representation engineering/steering, knowledge editing, unlearning, pruning) act on activations, parameters, or limited model structure, yet remain far cheaper and more targeted than full training. Some approaches also incorporate auxiliary components, such as compact classifiers or retrieval modules, that may require separate training; crucially, the core LLM weights remain unchanged.

To organize this diverse landscape, we propose a unified taxonomy of eight categories: \textbf{Context Engineering}, \textbf{Guardrails}, \textbf{Decoding Strategies}, \textbf{Representation Engineering}, \textbf{Editing}, \textbf{Unlearning}, \textbf{Pruning}, and \textbf{Multi-Agent Systems}. This taxonomy allows us to analyze inference-time methods not merely as ad-hoc fixes, but as a coherent design space for dynamic, modular, and multi-granular trustworthiness interventions.

% Addressing these multifaceted dimensions requires a diverse and sophisticated approach that can be applied dynamically and without costly retraining. This survey therefore concentrates on inference-time methods. We define these as techniques that guide, steer, or constrain LLM behavior during the generation process, without requiring full-scale retraining or fine-tuning of the entire model. These methods operate on a deployed model with significantly lower computational cost than training, encompassing everything from input-output filtering to targeted, localized modifications of the model's internal state or parameters. It also includes techniques that may leverage small, auxiliary models (e.g., classifiers or critics) that require their own training, so long as the primary LLM's weights remain frozen. The key principle is that these methods provide a dynamic and modular layer of control over a fixed, pre-trained model at the moment of use. Our investigation is organized around a comprehensive taxonomy of eight primary categories of such algorithms: Context Engineering, Guardrails, Decoding Strategies, Representation Engineering, Editing, Unlearning, Pruning, and Multi-Agent Systems. We structure these categories along a spectrum from external to internal interventions. External controls treat the model as an opaque black box, manipulating only its inputs and outputs; examples include prompt engineering and output filtering. Internal controls, in contrast, require white-box access to modify the model's parameters or internal states, such as directly altering activation vectors or editing model weights. This taxonomy allows us to logically group methods based on their level of intrusiveness and access requirements.

\begin{itemize}
    \item \textbf{Context Engineering:} Strategically designs the input prompt—rules, instructions, and few-shot exemplars (in-context learning)—to guide outputs without modifying the model.
    \item \textbf{Guardrails:} External safety modules that inspect inputs and/or outputs against predefined policies and then block, redact, or trigger regeneration when violations are detected.
    \item \textbf{Decoding Strategies:} Manipulate the token-level logit/probability distribution during generation to steer toward desired attributes (e.g., lower toxicity) or away from undesired ones.
    \item \textbf{Representation Engineering (Steering):} Fine-grained manipulation of internal activations at inference time. By adding or subtracting steering vectors tied to specific concepts (e.g., honesty), the method nudges the model’s computation in real time.
    \item \textbf{Editing:} Precise, localized parameter modifications that alter stored knowledge to correct facts, update information, or remove harmful associations at low computational cost.
    \item \textbf{Unlearning:} Selective removal of specific information, behaviors, or biases from a pre-trained model—important for data-forgetting requests and for excising harmful capabilities acquired during training.
    \item \textbf{Pruning:} Post-training removal of selected components (e.g., weights, neurons, attention heads). Traditionally for efficiency; its implications for trustworthiness are an emerging focus.
    \item \textbf{Multi-Agent Systems:} Replace single-model control with coordinated interactions among multiple LLM agents (e.g., debate, cross-verification), so that reliability emerges from structured collaboration.
\end{itemize}

% \begin{itemize}
% \item \textbf{Context Engineering:} This represents the most fundamental layer of control, leveraging the strategic design of the input provided to the model. It encompasses techniques such as prompt engineering, instruction crafting, and the use of few-shot examples (in-context learning) to guide the model's generative path and elicit desired responses without modifying the model itself.
% \item \textbf{Guardrails:} These are external modules, often rule-based or model-based, that function as safety filters. They operate either at pre-inference by examining user inputs or at post-inference by scrutinizing model outputs. Guardrails check for violations of predefined policies (e.g., safety, toxicity, privacy) and take corrective action, such as blocking, redacting, or regenerating content that fails to comply.
% \item \textbf{Decoding Strategies:} This category includes techniques that directly intervene in the token generation process. By manipulating the logit or probability distribution at each decoding step, these methods can steer the output towards desired attributes (e.g., lower toxicity, higher factuality) or embed signals (e.g., watermarking) without altering the model's parameters.
% \item \textbf{Representation Engineering (Steering):} A fine-grained approach that directly manipulates the model's internal state during inference. By intervening in the model's activation space—typically by adding or subtracting "steering vectors" corresponding to specific concepts (e.g., honesty, harmlessness)—this method can guide the model's computational process in real-time to align its outputs with desired values.
% \item \textbf{Editing:} These techniques aim to make precise, localized modifications to an LLM's parameters to alter its stored knowledge. Unlike full-scale retraining, model editing allows for surgically correcting factual inaccuracies, updating outdated information, or erasing specific memories, thereby enhancing factuality and privacy with minimal computational overhead.
% \item \textbf{Unlearning:} Conceptually the inverse of learning, unlearning methods focus on selectively removing specific information, behaviors, or biases from a pre-trained model. This is critical for addressing privacy concerns, such as fulfilling data forgetting requests under regulations like GDPR, and for mitigating safety risks by excising harmful capabilities acquired during training.
% \item \textbf{Pruning:} These methods involve the selective removal of model components—such as weights, neurons, or entire attention heads—after training is complete. While traditionally motivated by computational efficiency (i.e., reducing model size and latency), recent work has begun to explore the significant, often underexplored, secondary effects of pruning on trustworthiness, such as its potential to remove or amplify specific biases and behaviors.
% \item \textbf{Multi-Agent Systems:} This paradigm shifts control from a monolithic model to a collaborative system of multiple LLM agents. In this framework, trustworthiness is not enforced on a single agent but emerges from the architecture of their interactions. Mechanisms such as debate, cross-verification, and structured role-playing can produce more robust, factual, and nuanced outcomes than a single model could alone, as explored in recent work on bias detection and system robustness. 
% \end{itemize}

By synthesizing and critically evaluating these disparate lines of research, this survey aims to provide a unified conceptual framework for inference-time trustworthiness. We not only map the current landscape but also identify key synergies, unresolved challenges, and promising future directions at the intersection of these techniques. Our analysis reveals that a holistic approach, likely combining multiple methods from this taxonomy, will be essential for building the next generation of truly reliable and aligned AI systems.

\subsection{Related Surveys}
Our work is positioned at the confluence of three distinct but related research streams: inference efficiency, post-training capability enhancement, and broad-spectrum trustworthiness. By analyzing the scope and limitations of existing surveys in these areas, we can precisely define our unique contribution. 

% Our work builds upon and distinguishes itself from several existing lines of survey research. To clarify our unique contribution, we categorize related surveys into three primary areas: those focusing on inference efficiency, post-training modifications, and the broader topic of LLM trustworthiness. This structure allows us to pinpoint the specific gap our survey aims to fill.

\paragraph{Efficiency-focused Surveys.}
A significant body of work has surveyed inference-time optimizations for LLMs. For instance, comprehensive surveys like\citep{survey:inference_efficiency1} provide broad overviews of techniques such as quantization, knowledge distillation, and speculative decoding. Others focus on specific aspects, such as optimizations for Mixture-of-Experts models\citep{survey:inference_moe} or methods to reduce the memory footprint for deploying massive models\citep{survey:inference_efficiency2}. While foundational for understanding inference-time mechanics, their primary objective is computational performance (e.g., latency, throughput), not the explicit control of trustworthy behavior.

% A significant body of work has surveyed inference-time optimizations for LLMs. Several comprehensive surveys provide overviews of techniques to accelerate generation and reduce the computational cost of LLMs\citep{survey:inference_efficiency1, survey:inference_efficiency2, survey:inference_moe}. While foundational for understanding inference-time mechanics, their primary objective is computational performance (e.g., latency, throughput, memory).

\paragraph{Surveys on Post-Training Capability Enhancement.}
A second category of surveys examines methods applied after pre-training to improve a model's performance on complex tasks. These reviews often focus on enhancing emergent capabilities like reasoning or tool use, exploring advancements in reinforcement learning algorithms and prompting-based technics\citep{survey:post-training1, survey:post-training2}. The objective here is to make models *smarter* or more effective problem-solvers. While alignment is a peripheral concern, the central thrust is augmenting cognitive prowess, not enforcing the specific, granular constraints of safety, privacy, or fairness that are core to our work.

% A second category of surveys examines techniques applied after the main pre-training stage, such as methods for enhancing reasoning \citep{survey:post-training1, survey:post-training2}. However, their primary objective is fundamentally different from ours. These studies are largely concerned with enhancing the model's general capabilities and task-specific performance---essentially, making the model a more powerful reasoner or a more effective problem-solver. While alignment is sometimes discussed, the significant emphasis is on augmenting the model's intellectual prowess.

\paragraph{Trustworthy-focused Surveys.}
Finally, several recent surveys have addressed the crucial topic of LLM trustworthiness.
Broad, holistic surveys offer comprehensive overviews, charting the landscape of trustworthy LLMs by covering principles from data integrity to alignment evaluation\citep{survey:trustworthy1, survey:trustworthy2}.
While these works touch upon various mitigation strategies, they do not provide a systematic and unified taxonomy specifically for the post-hoc control methods we investigate.
Other research offers deep dives into specific facets or mechanisms of trustworthiness.
For instance, specialized surveys focus on mitigating particular harms like social bias in text-to-image generation\citep{survey:trustworthy_text2img}, ensuring safety in emerging applications like autonomous LLM agents\citep{survey:trustworthy_agents}, or exploring specific techniques such as machine unlearning to remove unwanted data or behaviors\citep{survey:unlearning1, survey:unlearning2}.
While invaluable, these works analyze individual challenges (e.g., bias, agent safety) or solution categories (e.g., unlearning) in isolation. They do not synthesize these different problems and techniques under the common, unifying framework of inference-time methods that form the basis of our work.

% While this body of work is foundational, a crucial gap remains. Prior surveys organize the field by the \textit{type of harm} (e.g., bias), a \textit{specific capability} (e.g., reasoning), or a \textit{single solution} (e.g., unlearning). To the best of our knowledge, no existing survey provides a systematic, cross-cutting analysis of the diverse techniques that constitute \textbf{inference-time methods for trustworthiness}. Our work fills this gap by proposing a unified taxonomy organized not by the problem, but by the mechanism of intervention—from external guardrails to internal representation steering. This reframing reveals the shared principles, unique trade-offs, and potential synergies between otherwise disparate methods, offering a structured guide for researchers and practitioners building the next generation of reliable AI systems.

\subsection{Contributions and Organization}
This survey provides a comprehensive analysis of inference-time methods for enhancing the trustworthiness of Large Language Models (LLMs). Our primary contribution is a novel, unified taxonomy that organizes the fragmented landscape of post-hoc control techniques into eight distinct categories, structured logically from external controls to internal manipulations: Context Engineering, Guardrails, Decoding Strategies, Representation Engineering, Editing, Unlearning, Pruning, and Multi-Agent Systems. We use this taxonomy as a lens to systematically evaluate how these methods address four core dimensions of trustworthiness: Safety, Privacy, Fairness, Factuality.

Our key contributions are as follows:

\begin{itemize}
    \item \textbf{A Unified Taxonomy of Inference-Time Methods for Trustworthiness.} We introduce a comprehensive taxonomy that organizes inference-time control methods into eight distinct categories. This framework structures the field by progressing from external controls (e.g., Context Engineering, Guardrails) to internal model manipulations (e.g., Representation Engineering, Model Editing) and systemic approaches (e.g., Multi-Agent Systems). This framework provides a structured lens through which to understand and compare how these methods can be leveraged to enhance LLM trustworthiness.
    \item \textbf{Systematic Cross-Dimensional Analysis.} We formulate trustworthiness as a multi-objective problem and conduct an in-depth analysis of how each method in our taxonomy addresses four core dimensions: Safety, Privacy, Fairness, and Factuality. This provides a holistic perspective on the capabilities and limitations of each technique, highlighting the inherent trade-offs between different trustworthiness goals.
    \item \textbf{Future Directions and Open Challenges.} We synthesize our findings to identify critical open challenges and articulate a forward-looking research agenda. We critically evaluate the scalability and practical limitations of current methods, discuss the pressing need for standardized evaluation benchmarks, and outline promising directions for building the next generation of robustly aligned AI systems.
\end{itemize}

% \begin{itemize}
%     \item \textbf{Propose a comprehensive, unified taxonomy} of eight distinct categories of low-cost, post-training control methods, bringing structure to a fragmented field.
%     \item \textbf{Provide a systematic, cross-dimensional analysis} of how these eight methods impact the core dimensions of trustworthiness (Safety, Privacy, Fairness, Factuality, and Faithfulness).
%     \item \textbf{Critically evaluate the inherent trade-offs} between different trustworthiness dimensions and performance metrics, while also assessing the adequacy of existing evaluation benchmarks.
%     \item \textbf{Identify key open challenges and articulate a forward-looking research agenda} to guide future work in this critical area.
% \end{itemize}

The remainder of the survey is organized as follows. \textbf{Sections \ref{sec:context_engineering}–\ref{sec:multi_agent_systems}} present the eight categories in our taxonomy and provide the cross-dimensional analysis, addressing our first and second contributions.. \textbf{Section \ref{sec:evaluation}} then provides a critical analysis of evaluation benchmarks, addressing the trade-offs and assessment challenges central to our third contribution. Finally, \textbf{Section \ref{sec:challenges}} consolidates open challenges and future directions. \textbf{Section \ref{sec:conclusion}} concludes the survey.
